% GENERATED FILE. Edit canonical lessons, then run npm run generate:books.
% canonical-source-hash: fnv1a64:ebe22da55d026219
% canonical-lessons: SA-C155-nutanah, SA-C155-jirnah, SA-C155-usnah, SA-C155-sitah, SA-C155-kathorah

\chapter{Fresh, Worn out, Hot, Chilly, Harsh}
\label{ch:sa-fresh-worn-out-hot-chilly-harsh}
\hlchaptermodality{155}

\begin{chapteropening}
\textbf{By the end of this chapter:} \emph{I can say fresh, worn out, hot, chilly and harsh.}

Complete the last lesson of chapter 155: I can say fresh, worn out, hot, chilly and harsh.
\end{chapteropening}

\section[\texorpdfstring{\sk{नूतनः} (nūtanaḥ)}{नूतनः (nūtanaḥ)}]{\sk{नूतनः} (nūtanaḥ) — fresh, new}

\label{lesson:SA-C155-nutanah}



\begin{quote}
Before the new one: say the Sanskrit for although, then the Sanskrit for or else.
\end{quote}

\subsection*{You'll want to know: \sk{नूतनः}}
\textbf{\sk{नूतनः}} — \emph{nūtanaḥ} — \textquotedblleft{}fresh, new\textquotedblright{}.

\begin{grammarlens}[title={What you've built}]
One more word for this chapter.
\end{grammarlens}

\subsection*{Your turn}
\begin{itemize}
\raggedright
  \item \emph{Say it:} \emph{nūtanaḥ}
  \item \emph{Say it:} \emph{nūtanaḥ}, once more
  \item \emph{Say it:} say \emph{nūtanaḥ}
  \item \emph{Recall:} say the Sanskrit for although, then the Sanskrit for or else, then say \emph{nūtanaḥ} again
\end{itemize}

\begin{tcolorbox}[breakable,colback=teal!4,colframe=teal!35!black,title={Before you move on}]
What does \sk{नूतनः} mean? (\textquotedblleft{}Fresh, new\textquotedblright{}.) Say it once more.
\end{tcolorbox}
\section[\texorpdfstring{\sk{जीर्णः} (jīrṇaḥ)}{जीर्णः (jīrṇaḥ)}]{\sk{जीर्णः} (jīrṇaḥ) — worn out, old}

\label{lesson:SA-C155-jirnah}



\begin{quote}
Before the new one: say the Sanskrit for or else, then the Sanskrit for fresh.
\end{quote}

\subsection*{You'll want to know: \sk{जीर्णः}}
\textbf{\sk{जीर्णः}} — \emph{jīrṇaḥ} — \textquotedblleft{}worn out, old\textquotedblright{}.

\begin{grammarlens}[title={What you've built}]
One more word for this chapter.
\end{grammarlens}

\subsection*{Your turn}
\begin{itemize}
\raggedright
  \item \emph{Say it:} \emph{jīrṇaḥ}
  \item \emph{Say it:} \emph{jīrṇaḥ}, once more
  \item \emph{Say it:} say \emph{jīrṇaḥ}
  \item \emph{Recall:} say the Sanskrit for or else, then the Sanskrit for fresh, then say \emph{jīrṇaḥ} again
\end{itemize}

\begin{tcolorbox}[breakable,colback=teal!4,colframe=teal!35!black,title={Before you move on}]
What does \sk{जीर्णः} mean? (\textquotedblleft{}Worn out, old\textquotedblright{}.) Say it once more.
\end{tcolorbox}
\section[\texorpdfstring{\sk{उष्णः} (uṣṇaḥ)}{उष्णः (uṣṇaḥ)}]{\sk{उष्णः} (uṣṇaḥ) — hot}

\label{lesson:SA-C155-usnah}



\begin{quote}
Before the new one: say the Sanskrit for fresh, then the Sanskrit for worn out.
\end{quote}

\subsection*{You'll want to know: \sk{उष्णः}}
\textbf{\sk{उष्णः}} — \emph{uṣṇaḥ} — \textquotedblleft{}hot\textquotedblright{}.

\begin{grammarlens}[title={What you've built}]
One more word for this chapter.
\end{grammarlens}

\subsection*{Your turn}
\begin{itemize}
\raggedright
  \item \emph{Say it:} \emph{uṣṇaḥ}
  \item \emph{Say it:} \emph{uṣṇaḥ}, once more
  \item \emph{Say it:} say \emph{uṣṇaḥ}
  \item \emph{Recall:} say the Sanskrit for fresh, then the Sanskrit for worn out, then say \emph{uṣṇaḥ} again
\end{itemize}

\begin{tcolorbox}[breakable,colback=teal!4,colframe=teal!35!black,title={Before you move on}]
What does \sk{उष्णः} mean? (\textquotedblleft{}Hot\textquotedblright{}.) Say it once more.
\end{tcolorbox}
\section[\texorpdfstring{\sk{शीतः} (śītaḥ)}{शीतः (śītaḥ)}]{\sk{शीतः} (śītaḥ) — chilly, cold}

\label{lesson:SA-C155-sitah}



\begin{quote}
Before the new one: say the Sanskrit for worn out, then the Sanskrit for hot.
\end{quote}

\subsection*{You'll want to know: \sk{शीतः}}
\textbf{\sk{शीतः}} — \emph{śītaḥ} — \textquotedblleft{}chilly, cold\textquotedblright{}.

\begin{grammarlens}[title={What you've built}]
One more word for this chapter.
\end{grammarlens}

\subsection*{Your turn}
\begin{itemize}
\raggedright
  \item \emph{Say it:} \emph{śītaḥ}
  \item \emph{Say it:} \emph{śītaḥ}, once more
  \item \emph{Say it:} say \emph{śītaḥ}
  \item \emph{Recall:} say the Sanskrit for worn out, then the Sanskrit for hot, then say \emph{śītaḥ} again
\end{itemize}

\begin{tcolorbox}[breakable,colback=teal!4,colframe=teal!35!black,title={Before you move on}]
What does \sk{शीतः} mean? (\textquotedblleft{}Chilly, cold\textquotedblright{}.) Say it once more.
\end{tcolorbox}
\section[\texorpdfstring{\sk{कठोरः} (kaṭhoraḥ)}{कठोरः (kaṭhoraḥ)}]{\sk{कठोरः} (kaṭhoraḥ) — harsh, rough}

\label{lesson:SA-C155-kathorah}



\begin{quote}
Before the new one: say the Sanskrit for hot, then the Sanskrit for chilly.
\end{quote}

\subsection*{You'll want to know: \sk{कठोरः}}
\textbf{\sk{कठोरः}} — \emph{kaṭhoraḥ} — \textquotedblleft{}harsh, rough\textquotedblright{}.

\begin{grammarlens}[title={What you've built}]
One more word for this chapter.
\end{grammarlens}

\subsection*{Your turn}
\begin{itemize}
\raggedright
  \item \emph{Say it:} \emph{kaṭhoraḥ}
  \item \emph{Say it:} \emph{kaṭhoraḥ}, once more
  \item \emph{Say it:} say \emph{kaṭhoraḥ}
  \item \emph{Recall:} say the Sanskrit for hot, then the Sanskrit for chilly, then say \emph{kaṭhoraḥ} again
\end{itemize}

\begin{tcolorbox}[breakable,colback=teal!4,colframe=teal!35!black,title={Before you move on}]
What does \sk{कठोरः} mean? (\textquotedblleft{}Harsh, rough\textquotedblright{}.) Say it once more.
\end{tcolorbox}
